% Draft for the Overleaf report, section "Data Transformation".
% Paste in place of the empty \textbf{a.} / \textbf{b.} in that section.
% Continues the existing numbering: this is Table 7 (Table 6 was the outlier
% comparison in Task 3; there is no new figure here).

\textbf{a.}
The only categorical columns are \texttt{Area} (200 countries) and \texttt{Item} (117 crops); \texttt{Year} is already numeric (Task 1), and \texttt{imputed} and the \texttt{*\_capped} flags from Tasks 2--3 are boolean indicators, not categories, so neither is encoded here. We use one-hot encoding: each unique value of \texttt{Area} and \texttt{Item} gets its own binary (0/1) column, e.g.\ \texttt{Area\_Afghanistan}, \texttt{Area\_Albania}, \ldots, \texttt{Item\_Almonds, with shell}, \ldots. This adds 317 new columns in total (200 for \texttt{Area}, 117 for \texttt{Item}).

Both \texttt{Area} and \texttt{Item} are nominal categories with no natural ordering -- one country or crop is not "more" or "less" than another. One-hot encoding is the appropriate choice for exactly this reason: it treats every category as equally distinct from every other, and introduces no rank between them. We deliberately avoided label encoding (assigning each category a single integer, e.g.\ Afghanistan~=~0, Albania~=~1, \ldots), because it would have implied a false ordering and a false distance between categories that do not exist in reality -- for instance, that Albania (1) is closer to Algeria (2) than to Brazil (20). Such an artefact can mislead models that interpret these values as distances or coefficients, and is particularly damaging for Task 6 (PCA): PCA is a linear technique that extracts components directly from the variance and covariance of its input columns, so a spurious alphabetical rank would be read as real structure and could distort the components. One-hot columns carry no such ordering, so any structure PCA finds among them reflects genuine patterns in the data rather than an artefact of the encoding -- which is the main reason we chose one-hot for this dataset.

We do not consider the resulting column count a problem. With 88,735 rows, 317 additional binary columns is a normal, manageable ratio of rows to features, and standard tooling represents one-hot columns as a sparse matrix, so the added columns cost little in memory or computation. The one setting where a high column count is a genuine concern is distance-based methods such as k-nearest neighbours, where many sparse binary dimensions can dilute the distance metric; since our downstream use is regression and PCA rather than k-NN, this is not a concern here.

\textbf{b.}
Task 3 already applied $\log_{10}(x+1)$ to area, production and yield and capped the remaining outliers per crop, correcting the strong right-skew (Task 1b) into distributions that are approximately symmetric and more "centred" than on the raw scale. Scaling addresses a separate problem: even after the transform, the three log-scaled measurements still span different ranges (roughly 0--7 for area, 0--9 for production, 0--7 for yield), and distance- or gradient-based models are sensitive to such differences in scale.

We use standardisation ($x' = (x-\bar{x})/s$) rather than min-max normalisation. Precisely because the outliers were already transformed in Task 3 and the values are therefore more centred, the mean and standard deviation are now more meaningful measures of spread than the min and max -- min-max would still be anchored to each column's two remaining endpoints, compressing the rest of the data near the middle of the $[0,1]$ range.

The scaling parameters are computed from the training set from Task 5 only (Table~\ref{tab:scaling}), and the same parameters are then used to transform both the training and the test set. In other words, we split the dataset before scaling it, not after. The reason is to avoid data leakage: had we instead computed the mean and standard deviation on the full dataset before splitting, or fitted a separate scaler per set, the test set's distribution would have influenced how the training data is represented, giving an artificially optimistic picture of how well the model generalises to unseen data. After scaling, the training set has mean 0 and standard deviation 1 by construction; the test set lands close to, but not exactly at, 0/1 (e.g.\ $-0.005/0.990$ for area), precisely because it is transformed with the training set's parameters rather than its own -- a check that the procedure was done correctly. The one-hot columns for \texttt{Area} and \texttt{Item} (Task 4a), \texttt{Year}, \texttt{imputed} and the \texttt{*\_capped} flags are not scaled, since they are binary indicators or discrete codes, not continuous measurements.

\begin{table}[H]
    \centering
    \begin{tabular}{lrr}
        \toprule
        Column & Mean & Std. dev. \\
        \midrule
        \texttt{area\_harvested\_ha\_log10} & 3.567 & 1.246 \\
        \texttt{production\_tonnes\_log10} & 4.283 & 1.310 \\
        \texttt{yield\_hg\_per\_ha\_log10} & 4.722 & 0.585 \\
        \bottomrule
    \end{tabular}
    \caption{Standardisation parameters for the log-scaled measurements, computed on the training set only.}
    \label{tab:scaling}
\end{table}
